\section{Formal exponentials and the complete word expansion}\label{sec:formal}
\subsection{The degree completion}
Let $V$ be the vector space over $\kk$ with basis $X,Y$, and put
\[
 T(V)=\bigoplus_{n\geq0}V^{\otimes n}=\kk\langle X,Y\rangle,
 \qquad
 \widehat T(V)=\prod_{n\geq0}V^{\otimes n}=\kk\langle\!\langle X,Y\rangle\!\rangle.
\]
The space $V^{\otimes n}=T_n$ has as basis the $2^n$ words of length $n$, the
empty word being the unit $1$. Multiplication concatenates words and is
extended to $\widehat T$ by the Cauchy product: the degree-$n$ part of a
product $AB$ is $\sum_{p+q=n}A_pB_q$, a finite sum. The filtration
$F^m\widehat T=\prod_{n\geq m}V^{\otimes n}$ satisfies
$F^mF^n\subseteq F^{m+n}$. A sum $\sum_j A_j$ of elements of $\widehat T$
\emph{converges formally} if for every degree $n$ only finitely many $A_j$
have a nonzero degree-$n$ part; its value is then defined degree by degree. No
norm is involved. In particular, if $A\in F^1\widehat T$, then $A^k\in
F^k\widehat T$, so for any scalar sequence $(c_k)$ the series
$\sum_{k\geq0}c_kA^k$ converges formally. Consequently substitution of
positive-degree series into any one-variable formal power series is well
defined, and
\begin{equation}\label{eq:explog}
 \exp A=\sum_{k\geq0}\frac{A^k}{k!},\qquad
 \log(1+U)=\sum_{k\geq1}\frac{(-1)^{k-1}}{k}U^k
 \qquad(A,U\in F^1\widehat T)
\end{equation}
are well-defined elements of $\widehat T$.

\begin{lemma}[Formal exponential and logarithm]\label{lem:explog}
The maps $\exp:F^1\widehat T\to1+F^1\widehat T$ and
$\log:1+F^1\widehat T\to F^1\widehat T$ of \eqref{eq:explog} are mutually
inverse bijections. If $A,B\in F^1\widehat T$ commute, then $e^{A+B}=e^Ae^B$.
If $g,h\in1+F^1\widehat T$ commute, then $\log(gh)=\log g+\log h$ and
$\log(g^{-1})=-\log g$.
\end{lemma}
\begin{proof}
We first prove the scalar identities $\log(e^t)=t$ and $e^{\log(1+t)}=1+t$
in $\Q[[t]]$. Put $E(t)=\sum_{k\geq0}t^k/k!$ and
$L(u)=\sum_{k\geq1}(-1)^{k-1}u^k/k$. Formal differentiation gives $E'=E$ and
$L'(u)=\sum_{k\geq1}(-1)^{k-1}u^{k-1}=(1+u)^{-1}$. The composition
$L(E(t)-1)$ has derivative $E'(t)/(1+E(t)-1)=E(t)/E(t)=1$ and vanishes at
$t=0$; a formal power series is determined by its derivative and constant
term, so $L(E(t)-1)=t$. Likewise $F(u)=E(L(u))$ satisfies
$F'(u)=E(L(u))L'(u)=F(u)/(1+u)$; the series $1+u$ satisfies the same
first-order equation with the same constant term $1$, and a solution of
$(1+u)F'=F$ is determined by its constant term because comparing the
coefficient of $u^n$ gives $(n+1)F_{n+1}=(1-n)F_n$ recursively. Hence
$E(L(u))=1+u$.

Now let $A\in F^1\widehat T$. The powers of the single element $A$ commute
with each other, so the subalgebra $\kk[[A]]$ of $\widehat T$ generated by $A$
is commutative, and the substitution $t\mapsto A$ defines a ring homomorphism
from $\Q[[t]]$ into $\widehat T$, well defined because each coefficient of
the image is a finite sum. Applying this homomorphism to $L(E(t)-1)=t$ gives
$\log(\exp A)=A$; applying the analogous homomorphism $u\mapsto U$ for
$U\in F^1\widehat T$ to $E(L(u))=1+u$ gives $\exp(\log(1+U))=1+U$. This
proves that the two maps are mutually inverse bijections.

If $A$ and $B$ commute, the binomial theorem holds for $(A+B)^n$, and
\[
 \sum_{n\geq0}\frac{(A+B)^n}{n!}
 =\sum_{n\geq0}\sum_{r+s=n}\frac{A^rB^s}{r!s!}
 =\sum_{r,s\geq0}\frac{A^rB^s}{r!s!}=e^Ae^B,
\]
all rearrangements being finite in each degree. If $g,h\in1+F^1\widehat T$
commute, then $\log g$ and $\log h$ commute, since each is a series in one
of the commuting elements $g-1$, $h-1$; hence
$e^{\log g+\log h}=e^{\log g}e^{\log h}=gh$, and applying $\log$ gives
$\log(gh)=\log g+\log h$. Taking $h=g^{-1}$, which exists as the geometric
series $\sum_{k\geq0}(1-g)^k$ and commutes with $g$, gives
$\log(g^{-1})=-\log g$. No commutativity of unrelated elements was used.
\end{proof}

In a Banach algebra the exponential series converges absolutely for every
$A$, since $\norm{A^n/n!}\leq\norm A^n/n!$, and the logarithm series converges
absolutely whenever $\norm U<1$. The scalar identities of the proof then hold
in the closed commutative subalgebra generated by $A$ or $U$, by the same
substitution argument with absolutely convergent series, so the analytic
statements corresponding to \cref{lem:explog} hold on those domains.

\subsection{Definition of the BCH series and the block expansion}
Define
\begin{equation}\label{eq:Zdef}
 Z(X,Y)=\BCH(X,Y)=\log(e^Xe^Y)=\sum_{n\geq1}Z_n(X,Y),\qquad Z_n\in V^{\otimes n}.
\end{equation}
By \cref{lem:explog}, $Z$ is the unique element of $F^1\widehat T$ with
$e^Z=e^Xe^Y$; thus \eqref{eq:main} already holds as an associative formal
identity. Expanding the two exponentials first and then the logarithm gives
\begin{equation}\label{eq:blocks}
 Z(X,Y)=\sum_{k\geq1}\frac{(-1)^{k-1}}{k}
 \sum_{\substack{r_i,s_i\geq0\\r_i+s_i>0\ (1\leq i\leq k)}}
 \frac{X^{r_1}Y^{s_1}\cdots X^{r_k}Y^{s_k}}
 {\prod_{i=1}^k r_i!s_i!}.
\end{equation}
Indeed $e^Xe^Y-1=\sum_{r+s>0}X^rY^s/(r!s!)$, and the $k$th power of this sum
is the sum over all $k$-tuples of nonempty blocks $X^{r_i}Y^{s_i}$. This
proves the associative expansion displayed on the source page. At this stage
the inner sum is noncommutative multiplication, not a sum over independent
commutators, and nothing has yet been proved about the Lie nature of the
homogeneous parts.

\subsection{A finite nonrecursive formula for every word}
For a nonempty word $w$ define its \emph{block weight}
\begin{equation}\label{eq:weight}
 a(w)=
 \begin{cases}
 \dfrac1{r!s!},&w=X^rY^s,\quad r,s\geq0,\quad r+s>0,\\[6pt]
 0,&\text{otherwise}.
 \end{cases}
\end{equation}
Thus $a(w)=0$ precisely when $w$ contains the adjacent pattern $YX$; pure
powers of one letter are permitted. A \emph{cut} of $w$ into $k$ pieces is a
factorization $w=w_1\cdots w_k$ into consecutive nonempty subwords.

\begin{theorem}[Every associative BCH coefficient]\label{thm:wordcut}
For every nonempty word $w=a_1\cdots a_n$,
\begin{equation}\label{eq:wordcut}
 c(w):=\coeff{w}Z(X,Y)
 =\sum_{k=1}^{n}\frac{(-1)^{k-1}}{k}
   \sum_{\substack{w=w_1\cdots w_k\\|w_i|>0}}\prod_{i=1}^k a(w_i)
 =\sum_{k=1}^{n}\frac{(-1)^{k-1}}{k}
   \sum_{0=i_0<i_1<\cdots<i_k=n}
   \prod_{j=1}^{k}a\bigl(a_{i_{j-1}+1}\cdots a_{i_j}\bigr).
\end{equation}
Consequently
\begin{equation}\label{eq:homblocks}
 Z_n=\sum_{w\in\{X,Y\}^n}c(w)\,w
 =\sum_{k=1}^{n}\frac{(-1)^{k-1}}{k}
 \sum_{\substack{r_i+s_i>0\\\sum_i(r_i+s_i)=n}}
 \frac{X^{r_1}Y^{s_1}\cdots X^{r_k}Y^{s_k}}
 {\prod_i r_i!s_i!}.
\end{equation}
Both formulas are finite and nonrecursive, and every coefficient is rational.
For a word of length $n$ there are exactly $2^{n-1}$ cut patterns before zero
terms are discarded.
\end{theorem}
\begin{proof}
Write $U=e^Xe^Y-1\in F^1\widehat T$. Its coefficient on a nonempty word $v$
is $a(v)$: the only words with nonzero coefficient in $e^Xe^Y$ are the
products $X^rY^s$, which arise once each, with coefficient $1/(r!s!)$. The
coefficient of $w$ in $U^k$ is obtained by expanding the product of $k$
copies of $U$: a word appears in the product exactly when it is the
concatenation $w_1\cdots w_k$ of one word from each factor, and its
coefficient is the sum over all such factorizations of the product of the
factor coefficients. Since every word in $U$ is nonempty, these
factorizations are precisely the cuts of $w$ into $k$ nonempty pieces, and
the coefficient is $\sum_{w=w_1\cdots w_k}\prod_ia(w_i)$. Because $w$ has
length $n$, no cut into more than $n$ nonempty pieces exists, so $U^k$ has
zero coefficient on $w$ for $k>n$. Substituting into
$Z=\sum_{k\geq1}(-1)^{k-1}U^k/k$ gives \eqref{eq:wordcut}; the second form
just parametrizes a cut by the positions $i_1<\cdots<i_{k-1}$ at which it
breaks. Choosing a subset of the $n-1$ gaps between consecutive letters is
the same as choosing a cut pattern, which gives the count $2^{n-1}$.
Expanding each permitted block as $X^{r_i}Y^{s_i}$ in \eqref{eq:wordcut}
and summing over words of length $n$ gives \eqref{eq:homblocks}, which is
the degree-$n$ part of \eqref{eq:blocks}.
\end{proof}

\begin{example}
For $w=XY$ the one-piece contribution is $a(XY)=1$ and the two-piece
contribution is $-\tfrac12a(X)a(Y)=-\tfrac12$, so $c(XY)=\tfrac12$. For
$w=YX$ the one-piece contribution is $a(YX)=0$, so $c(YX)=-\tfrac12$. For
$w=XXY$ the one-piece contribution is $a(XXY)=\tfrac12$; the two cuts
$X\mid XY$ and $XX\mid Y$ contribute $-\tfrac12(1\cdot1+\tfrac12\cdot1)
=-\tfrac34$; the three-piece cut contributes $\tfrac13$. Their sum is
$\tfrac12-\tfrac34+\tfrac13=\tfrac1{12}$. For $w=XYX$ the one-piece
contribution is zero, the two cuts $X\mid YX$ and $XY\mid X$ contribute
$-\tfrac12(0+1)=-\tfrac12$, and the three-piece cut contributes $\tfrac13$,
giving $c(XYX)=-\tfrac16$. These small computations already show why the
coefficients are rational and why cancellation between different cut
patterns matters.
\end{example}

\subsection{Any finite number of exponential factors}\label{sec:manyexp}
The same construction works without change for several factors.
\begin{corollary}[Several exponentials]\label{cor:manyexp}
Let $A_1,\ldots,A_m$ be letters. For a $k\times m$ array
$\rho=(r_{ji})$ of nonnegative integers in which every row has positive
sum, let $w(\rho)$ be the word obtained by concatenating the row words
$A_1^{r_{j1}}\cdots A_m^{r_{jm}}$ for $j=1,\ldots,k$, and put
$|\rho|=\sum_{j,i}r_{ji}$. Then
\begin{equation}\label{eq:manyexp}
 \log(e^{A_1}\cdots e^{A_m})
 =\sum_{k\geq1}\frac{(-1)^{k-1}}{k}
   \sum_{\rho}\frac{w(\rho)}{\prod_{j,i}r_{ji}!},
\end{equation}
and the coefficient of a given word $w$ is the finite sum over cuts of $w$
into $k$ nonempty pieces of the product of the block weights
$a_m(w_i)$, where $a_m(A_1^{r_1}\cdots A_m^{r_m})=(r_1!\cdots r_m!)^{-1}$
and $a_m$ vanishes on every other word.
Equivalently, with $t$ a formal scalar,
\begin{equation}\label{eq:multibch}
 \log\Bigl(\prod_{i=1}^{m}e^{tA_i}\Bigr)
 =\sum_{k\geq1}\frac{(-1)^{k-1}}{k}
 \sum_{\rho}\frac{t^{|\rho|}}{\prod_{j,i}r_{ji}!}\,w(\rho).
\end{equation}
\end{corollary}
\begin{proof}
Repeat the proof of \cref{thm:wordcut} with $U=e^{A_1}\cdots e^{A_m}-1$,
whose nonempty words are exactly the products $A_1^{r_1}\cdots A_m^{r_m}$
with coefficient $(r_1!\cdots r_m!)^{-1}$. The scalar $t$ simply counts
total degree.
\end{proof}
If the factors are of the form $e^{\alpha_jX}$ or $e^{\beta_jY}$ with the
letters $X,Y$ repeated in different positions, one first uses a distinct
letter for each factor and then substitutes; the resulting coefficients are
finite polynomials in the scalar weights $\alpha_j,\beta_j$. After
\cref{sec:hopf} it will follow that each homogeneous part of
\eqref{eq:manyexp} is a Lie polynomial and that each word $w(\rho)$ may be
replaced by $R(w(\rho))/|\rho|$. This observation supplies every order, not
only the leading error, for the splitting formulas of \cref{sec:splitting}.
